\chapter{三个调控环节：持续振荡与短暂适应}
双节点已经提供两条路线：正回路可以形成开关，负回路配上正自反馈则有二维振荡的入口。若没有正自边，前面的瞬时产生--一阶降解模型仍排除了周期轨道。第三个环节提供另一种时间结构：在负闭环中，它增加动态环节，让没有正自边的模型也能出现摆动失稳；在前馈中，它让同一输入沿快慢两条路径到达输出。下面分别讨论振荡、时间筛选与适应。
\begin{figure}[H]\centering
\begin{tikzpicture}
\node[gene] (a) at (0,1.4) {$X_1$};\node[gene] (b) at (3,1.4) {$X_2$};\node[gene] (c) at (1.5,-.5) {$X_3$};
\draw[rep] (a)--(b);\draw[rep] (b)--(c);\draw[rep] (c)--(a);
\node[gene] (u) at (6,1.4) {$u$};\node[gene] (y) at (9,1.4) {$y$};\node[gene] (z) at (7.5,-.5) {$z$};
\draw[act] (u)--(y);\draw[act] (u)--node[left]{慢支路}(z);\draw[rep] (z)--(y);
\end{tikzpicture}
\caption{左：三次抑制构成负反馈环。右：不一致前馈回路，直接激活与间接抑制在输出处竞争；输入 $u$ 在这里是外部信号。}
\end{figure}

\section{闭合负反馈：三个环节怎样使稳态失稳？}
三个相互抑制环节可写为（本节取相同的一阶降解率 $\gamma$）
\begin{equation}\dot x_1=f_1(x_3)-\gamma x_1,\quad\dot x_2=f_2(x_1)-\gamma x_2,\quad\dot x_3=f_3(x_2)-\gamma x_3.\end{equation}
反馈符号仍然是负，但稳定性已经不能照搬双节点的结论。要找出新的可能，先用稳态附近的斜率近似产生率曲线，计算三种基本偏差各自会增长还是衰减。记 $a_i=f_i'$ 为三个交叉斜率。三次抑制使 $a_1a_2a_3<0$，写成 $a_1a_2a_3=-k^3$、$k>0$，则
\[
\lambda I-J=\begin{pmatrix}\lambda+\gamma&0&-a_1\\-a_2&\lambda+\gamma&0\\0&-a_3&\lambda+\gamma\end{pmatrix},
\quad \det(\lambda I-J)=(\lambda+\gamma)^3-a_1a_2a_3.
\]
当乘积为 $-k^3$ 时，取 $-1$ 的三个立方根 $-1,\frac12\pm i\frac{\sqrt3}{2}$ 得到
\[
\lambda_1=-\gamma-k,\qquad\lambda_{2,3}=-\gamma+\frac k2\pm i\frac{\sqrt3k}{2}.
\]
$k<2\gamma$ 时，三种偏差的增长率都为负，会逐渐减小；$k>2\gamma$ 时，其中两种带摆动的偏差开始增长。门槛 $k=2\gamma$ 处，摆动的角频率为 $\sqrt3\gamma$，对应周期 $2\pi/(\sqrt3\gamma)$。
对对称 Hill 环 $f(s)=\alpha/(1+s^n)$，唯一对称稳态满足 $\alpha=\gamma s(1+s^n)$，因此
\[
\frac k\gamma=\frac{ns^n}{1+s^n}.
\]
这里也不会漏掉另外的不对称稳态：当三个调控函数相同且都严格下降时，沿环走过三次调控，浓度增加最终对应的是反馈输入下降。这条反馈关系与 $x$ 本身只会相交一次，而三个浓度相等的解已经满足关系，因此它就是整个环的唯一稳态。
当 $n\leq2$ 时不可能达到有限参数的阈值；当 $n>2$ 时
\[
s_c=\left(\frac2{n-2}\right)^{1/n},\qquad
\alpha_c=\gamma\frac n{n-2}\left(\frac2{n-2}\right)^{1/n}.
\]
在 $n=3,\gamma=1$ 下 $\alpha_c=3\sqrt[3]2$，演示的 $\alpha=20$ 位于稳态失稳一侧。
随 $\alpha$ 增加，$\dd s/\dd\alpha=1/[\gamma(1+(n+1)s^n)]>0$，且 $\dd k/\dd s>0$；所以随参数增大，这两种摆动方式确实会从衰减转为增长。这对复特征值的实部从负变正，数学上称为\emph{穿过虚轴}；“虚轴”就是复数平面中实部为零的位置。因为这里增长率随参数变化确实不为零，这次穿越是\emph{横截}的，也就是实际穿过去，而不是碰到边界后折返。这是 Hopf 分岔的候选点：稳态附近原本衰减的摆动可能变成周期振荡。

要进一步用标准 Hopf 分岔定理判断附近的周期解，还需模型足够光滑，并检验非线性的非退化条件。常用的量是\emph{第一 Lyapunov 系数}，它描述微小振幅开始增长后，非线性会限制这种增长还是进一步增强它。标准的非退化条件要求这个系数不为零；它的符号还帮助判断新出现的周期解是否稳定。这个系数在本文尚未计算，所以不能据此宣称已经证明稳定的周期振荡。以上计算只确定了小偏差不再衰减的门槛。图~\ref{fig:rep}采用 $f(s)=20/(1+s^3)$ 展示给定初值下的持续振荡；周期轨道的存在与稳定性属于进一步的非线性问题，不能由这幅图替代证明。


\begin{figure}[H]\centering\includegraphics[width=.78\linewidth]{figures/repressilator.pdf}\caption{三节点负反馈环的数值轨迹，生成脚本见 \texttt{code/make\_figures.py}。}\label{fig:rep}\end{figure}

\section{快慢前馈：输出怎样先响应再回落？}
输入 $u$ 可以直接激活输出 $y$，又通过慢变量 $z$ 抑制 $y$：
\begin{equation}\tau_y\dot y=H(u,z)-y,\qquad\tau_z\dot z=G(u)-z.\end{equation}
在 $H_u>0,H_z<0,G'>0$ 的调控范围内，若 $z$ 比 $y$ 慢，上升阶跃可使 $y$ 先增加，随后 $z$ 增加抑制 $y$，形成脉冲。\emph{脉冲}与\emph{精确适应}不同：精确适应要求稳态输出 $y^*(u)$ 与持续输入的幅度无关，这对 $H$、$G$ 有额外条件；拓扑本身不保证这种性质。后面的 Cell 论文讨论会把这一支路抵消机制与带缓冲节点的负反馈作具体比较。

要把“回落”与“回到同一基线”区分开，先算恒定输入下的长期输出。设 $\tau_y,\tau_z>0$，$H,G$ 在所考察输入及状态范围内连续，固定 $u$ 时有
\[
z(t)=G(u)+(z_0-G(u))e^{-t/\tau_z},
\]
所以 $z\to G(u)$。此时输入输出支路 $H(u,z(t))\to H(u,G(u))$，按前面单向调控的方法把 $y$ 的方程积分，就得到
\[
y^*(u)=H(u,G(u)).
\]
精确适应要求输入改变后仍回到相同的输出，因此恰好要求这个复合函数在输入区间内为常数。如果 $H,G$ 为 $C^1$，且所讨论的输入覆盖一个完整区间 $I$，那么可以用\emph{链式法则}沿两条支路计算输入变化的影响。链式法则就是：间接支路的总变化率等于沿路各步变化率的乘积。由此得到
\[
H_u(u,G(u))+H_z(u,G(u))G'(u)=0\qquad(u\in I).
\]
这个式子揭示了适应所需的抵消：直接输入对输出的长期增益，与输入通过 $z$ 产生的长期增益大小相等、符号相反。快慢差异帮助形成瞬态峰，但只有这一稳态抵消才使基线与输入幅度无关。

可以选一个简单函数把这两个层面同时展示出来。取 $u>0,z>0$、$G(u)=u$、$H(u,z)=u/z$。每个恒定输入的稳态为 $(y^*,z^*)=(1,u)$，故确有精确适应。若从输入 $u_0$ 的稳态出发，在 $t=0$ 切换至 $u_1>u_0$，则
\[
z(t)=u_1+(u_0-u_1)e^{-t/\tau_z},\qquad
\dot y(0^+)=\frac{u_1/u_0-1}{\tau_y}>0.
\]
$y$ 起初升高，最终回到 $1$，所以至少存在一个高于基线的峰；$\tau_z\gg\tau_y$ 使脉冲更易分辨。$u/z$ 是理想化的比值调控，只在 $u,z>0$ 的工作域内使用。若换成 $H=u/(1+z)$，保持同一拓扑与 $G$，则 $y^*=u/(1+u)$ 依赖输入，精确适应消失。这给出了拓扑不充分的明确反例。
\begin{figure}[H]\centering\includegraphics[width=.85\linewidth]{figures/feedforward_adaptation.pdf}
\caption{相同的不一致前馈结构可以精确适应，也可以只产生瞬态峰而改变最终基线。输入从 $1$ 跳至 $2$，$\tau_y=0.25,\tau_z=2$；右图比较两种调控函数的稳态输入输出曲线。}\label{fig:adapt}
\end{figure}

\section{一致前馈也值得研究：它可以筛选信号持续时间}
前馈回路有两条从输入到输出的路径。若两条路径的累计符号相同，称为\emph{一致前馈}；相反则为\emph{不一致前馈}。前面的例子利用两条路径相反的长期作用形成适应，而一致前馈的一种用途是区分短脉冲与持续输入。

看一个全激活的三节点回路 $A\to B\to C$、$A\to C$。为单独看时间逻辑，暂时把输入节点 $A$ 的活动当成 $0$ 或 $1$，令 $\tau_B\dot B=A-B$，再以阈值 $0<\theta<1$ 判断 $B$ 是否足够高。如果输出产生采用 AND 逻辑，要求 $A$ 与 $B$ 同时高：上升输入到来后，$B(t)=1-e^{-t/\tau_B}$，只有等到
\[
t_{\rm on}=-\tau_B\ln(1-\theta)
\]
才开始产生输出；输入一撤去，AND 条件立即失效。所以比这个时间还短的输入可能被过滤掉。如果改成 OR 逻辑，输入到来时直接支路立即打开，撤去后则要等 $B$ 从高值降到阈值，关闭的等待时间为 $t_{\rm off}=\tau_B\ln(1/\theta)$。

这是理想阈值模型，真实的连续产生率会把“立即”与“等待”平滑成不同快慢的响应，输出自身的恢复时间也会再加一层滞后。重点是：同样的三条激活边，配上不同启动子逻辑，就可能分别延迟开启或延迟关闭。因此一致前馈并非“没有适应所以无趣”，它回答的是另一个时间功能问题。

\section{汤超参与的 Cell 工作：从“所有图”转向“目标功能”}
Ma、Trusina、El-Samad、Lim 与汤超的 \emph{Defining Network Topologies that Can Achieve Biochemical Adaptation}（Cell, 2009，参考文献~6）提供了一种研究思路：先定功能标准，再在统一动力学规则下枚举。研究节点是总量固定、在活性与非活性状态之间转换的酶；连边按 Michaelis--Menten 反应建模，和本文的转录产生--降解模型不同。

三节点的九条潜在有向边给出 $3^9$ 种符号图。当输入固定在 $A$、输出固定在 $C$ 时，没有因果路径的图必须缺少 $A\to C$，而且 $A\to B$、$B\to C$ 不能同时存在。这两条边共有 $5$ 种不形成路径的组合，其余六条边任取，故筛去 $5\times3^6$ 种，剩下 $3^9-5\times3^6=16\,038$ 种。论文对这些拓扑各采样 $10\,000$ 组参数，用灵敏度与适应精度筛选。每图采样中超过 $10$ 组参数达标的 $395$ 个较稳健拓扑均含缓冲节点负反馈（NFBLB）或比例节点不一致前馈（IFFLP）核心。这里的数目是该模型与筛选规则的结果，不是所有三基因系统的功能计数。

为了理解筛选到底在看什么，设输入从 $I_1$ 变为 $I_2$，输出从初始稳态 $O_1$ 出发，达到最明显偏离的位置 $O_{\rm peak}$，最后趋于 $O_2$。取非零的基线和输入变化，定义
\[
\Delta_I=\frac{|I_2-I_1|}{|I_1|},\qquad
S=\frac{|O_{\rm peak}-O_1|/|O_1|}{\Delta_I},\qquad
E=\frac{|O_2-O_1|/|O_1|}{\Delta_I},\qquad P=E^{-1}.
\]
$S$ 是\emph{灵敏度}，问输出有没有明显响应；$P$ 是\emph{适应精度}，问最终偏差是否很小。论文采用 $S>1,P>10$ 的功能门槛。精确适应对应 $E=0$，形式上 $P$ 无限大，但这不是数值采样直接证明的性质。

一个始终不动的输出可以有很小的 $E$，却不是良好的变化探测器；一个不断振荡的输出也不能简单用最终平台 $O_2$ 评价。论文补充材料排除了持续振荡，并对衰减摆动设置了额外规则。对“适应”的筛选失败不等于对“振荡”的研究失败：在前面振荡例子中有用的时钟，在这里恰恰可能被排除。

\section{缓冲节点负反馈：把偏差积累起来，直到输出恢复}
为什么普通的负反馈不自动给出精确适应？因为一个通常的恢复变量会同时遗忘过去的偏差：产生与降解平衡时，允许输出停在一个小的残余误差上。要把这个误差消掉，调节变量需要在误差存在时继续改变。

下面用本文的简化三变量模型说明 NFBLB 的这一作用。$A$ 接收输入，$C$ 是输出，$B$ 是调节变量；$C$ 激活 $B$，$B$ 抑制 $C$。令
\begin{align}
\dot A&=\rho\{u(1-A)-A\},\\
\dot B&=\kappa(C-C_0),\\
\dot C&=k\{A(1-C)-BC\},\qquad \rho,\kappa,k>0.\label{eq:buffer_demo}
\end{align}
$A,C$ 可以理解为活性比例；$B$ 是缓冲变量，在本例的工作范围内也保持正且小于 $1$。第一式使 $A^*=u/(1+u)$。关键在第二式：当 $C>C_0$，$B$ 继续积累并压低 $C$；当 $C<C_0$，$B$ 下降，抑制减弱；只有 $C=C_0$ 时，$B$ 才停止调整。因此
\[
C^*=C_0,\qquad B^*=\frac{1-C_0}{C_0}A^*.
\]
输入变大后，长期改变的是 $A^*$ 与 $B^*$，输出仍为 $C_0$。这不是两个支路恰巧抵消一次，而是反馈持续校正误差。

把第二式积分，得到
\[
B(t)=B(0)+\kappa\int_0^t[C(s)-C_0]\dd s.
\]
所以 $B$ 保存的是过去输出误差的累积，称为\emph{积分反馈}。这里的“缓冲”不只是响应慢，而是 $\dot B$ 在工作近似下不直接依赖 $B$ 自己。饱和的酶反应能提供近似恒定的转化速率，从而产生这类条件；它不会无条件地出现在所有负反馈回路中。式~\eqref{eq:buffer_demo}是机制模型，不是论文全套速率方程；若 $B$ 达到上下限，线性积分近似就失效，必须回到完整的转化反应。

这个例子的恢复也能算出来。因为 $A$ 不受下游影响，线性化后它的一根特征值为 $-\rho(1+u)$，而 $(B,C)$ 部分为
\[
J_{BC}=\begin{pmatrix}0&\kappa\\-kC_0&-k(A^*+B^*)\end{pmatrix}.
\]
迹为负，行列式 $\kappa kC_0>0$，故只要 $0<C_0<1$ 且 $B^*$ 在缓冲近似的内部工作范围，稳态就局部渐近稳定。输出基线固定与基线附近能恢复，是这里分别检查过的两件事。

若给缓冲变量加入“泄漏”，改成
\[
\dot B=\kappa(C-C_0)-\varepsilon_B(B-B_{\rm ref}),\qquad\varepsilon_B>0,
\]
稳态就允许 $C^*-C_0=(\varepsilon_B/\kappa)(B^*-B_{\rm ref})$。它仍是负反馈，也仍可能有瞬态峰，却通常留下随输入变化的基线偏差。缓冲变量慢并不足够，是否能持续积累误差才是区别。

\section{比例节点前馈：长期随输入增加，恰好抵消直接支路}
前馈的思路不同：没有输出回传的校正，而是让一条支路产生与输入成比例的反向作用。沿用同样的输入节点 $A$，构造另一个简化模型：
\begin{align}
\dot A&=\rho\{u(1-A)-A\},\\
\dot B&=k_BA-\gamma_BB,\\
\dot C&=k_AA(1-C)-k_CB C,\qquad k_B,\gamma_B,k_A,k_C>0.
\end{align}
$A$ 直接激活 $C$，又通过 $B$ 抑制 $C$。令 $\eta=k_B/\gamma_B$，稳态下
\[
B^*=\eta A^*,\qquad
0=A^*\{k_A(1-C^*)-k_C\eta C^*\}.
\]
在 $A^*>0$ 的工作范围内，公共因子 $A^*$ 可以约去，得到
\[
C^*=\frac{k_A}{k_A+k_C\eta}.
\]
这就是比例节点的作用：输入增加多少，抑制支路长期也按固定比例增加多少，两条支路在输出方程中共享输入因子。$B$ 较慢时，直接支路先起作用，随后比例支路跟上，于是输出先偏离、再回到基线。

该模型是逐级驱动的：$A$ 先稳定，随后 $B$ 稳定，最后 $C$ 的恢复率为 $k_AA^*+k_CB^*>0$，所以在给定正输入及上述工作范围内三者都收敛。若把 $B$ 的产生换成 $k_BA/(K+A)$，则 $B^*/A^*=k_B/[\gamma_B(K+A^*)]$ 不再恒定，输出基线也通常随输入改变。可见不一致前馈的边符号还要配上合适的支路比例关系。

\begin{figure}[H]\centering
\begin{tikzpicture}
\node[gene] (a) at (0,1.3) {$A$};\node[gene] (c) at (3,1.3) {$C$};\node[gene] (b) at (3,-.4) {$B$};
\draw[act] (a)--(c);\draw[act] (c) to[bend left=22] (b);\draw[rep] (b) to[bend left=22] (c);
\node at (1.5,-1.15) {缓冲：$\dot B\propto C-C_0$};
\node[gene] (aa) at (6,1.3) {$A$};\node[gene] (cc) at (9,1.3) {$C$};\node[gene] (bb) at (7.5,-.4) {$B$};
\draw[act] (aa)--(cc);\draw[act] (aa)--(bb);\draw[rep] (bb)--(cc);
\node at (7.5,-1.15) {比例：$B^*\propto A^*$};
\end{tikzpicture}
\caption{两种适应机制的节点角色。左边由输出误差驱动缓冲变量持续校正；右边由输入驱动比例变量，长期抵消直接输入作用。两者都允许 $B$ 的稳态随输入改变。}
\end{figure}
\begin{figure}[H]\centering\includegraphics[width=\linewidth]{figures/adaptation_mechanisms.pdf}
\caption{本文简化模型的数值比较。左：缓冲反馈与比例前馈都先响应再恢复，带泄漏的缓冲反馈改变最终基线。中：调节变量 $B$ 保留输入变化。右：理想机制的稳态输出与输入无关，泄漏模型则不然。曲线按各自初始基线归一化；并非复现论文的全网参数扫描。}\label{fig:adaptmech}
\end{figure}
本图的示例参数为 $\rho=3$、输入 $u:0.5\to0.8$。缓冲模型取 $C_0=0.4,\kappa=0.3,k=2$，泄漏模型另取 $\varepsilon_B=0.08$；比例模型取 $k_B=\gamma_B=0.25,k_A=k_C=2$。这些参数用于展示基线恢复与泄漏反例，没有按论文的 $S>1$ 灵敏度门槛选取，不能把本图直接计入作者的功能筛选结果。完整方程与检查见 \texttt{code/extended\_examples.py}。

\section{为什么功能枚举不能只看节点数与连边？}
论文的补充材料进一步显示：同一个负反馈环，把哪个节点当作输出会改变是否适应；带正自边的特定速率条件还可让 $\log B$ 累积输出误差。例如补充材料的一种近似为 $\dot B=\nu B(C_0-C)$，在 $B>0$ 时可写成
\[
\frac{\dd\log B}{\dd t}=\nu(C_0-C).
\]
于是存储累计误差的是 $\log B$，内部稳态仍由 $C=C_0$ 决定。这个例子涉及饱和速率条件及随调控变量改变的损失作用，不能直接套用固定一阶降解模型的分类结论。因此“缓冲节点”与“比例节点”是动态角色，不只是绘图中的某个位置。两个例子也提醒我们，正自反馈既可能存储状态，也可能参与适应；功能要由方程中的作用方式来区分（参考文献~7）。

论文还扫描了固定输入节点 $A$、输出节点 $C$ 的 $81$ 种双节点酶网络，在所扫参数范围内没有得到符合其标准的适应。它没有证明任意两变量系统都不能适应：比如把输入直接施加于输出，再用另一个变量积分误差，就改变了输入、输出角色；换用不同降解机制也改变了模型族。本文的两节点分类讨论的是恢复、记忆和 Hopf 入口，而论文在同样数量的结构上问的是另一个有指定标准的功能问题。

“稳健”也要和“精确”区分。精确适应问某个模型的稳态输出是否在输入区间内恒定；稳健性问参数不必调到孤立的特殊数值时，这个功能是否仍能实现。数值上可以记录同一拓扑的 $N$ 组参数中有多少组符合标准，写成 $Q$ 或比例 $Q/N$。$Q$ 依赖参数范围、采样规则、输入步骤和功能门槛，不能当作拓扑唯一的内在属性。它也不要求输出基线对所有参数变化都不变：同一组参数下适应不同输入，与换参数后保持同一个基线，是两个不同要求。

所以更彻底的三节点研究应同时记录\emph{拓扑、节点角色、速率机制、参数、功能标准}。一致前馈可以筛选信号持续时间，不一致前馈可以形成脉冲并在比例条件下适应，负反馈既可能恢复或振荡，也可能在缓冲条件下适应。对一项功能不合格的网络，可能恰好实现另一项功能；这比给所有图固定贴上“有趣”或“无趣”的标签更有解释力。
